% Generated by docs/generate_user_guide.py. Do not edit directly.
% Source: utilities/descriptions.yaml: cv_schemes

\begin{longtable}{@{}>{\raggedright\arraybackslash}p{0.20\textwidth}>{\raggedright\arraybackslash}p{0.28\textwidth}>{\raggedright\arraybackslash}p{0.44\textwidth}@{}}
\toprule
\textbf{Scheme} & \textbf{Internal key} & \textbf{Description} \\
\midrule
\endfirsthead
\toprule
\textbf{Scheme} & \textbf{Internal key} & \textbf{Description} \\
\midrule
\endhead
Shuffle split & shuffle\_split & Repeatedly creates random training and test subsets. This is the default because it provides a general robustness check of whether the selected curve form captures the overall response pattern. It is not time-aware, so random splits may make results look more reassuring when nearby observations are strongly related. The splits, test-size, and random-seed settings control the repeated random splits. \\
Leave-one-timepoint-out & leave\_one\_timepoint\_out & Sets aside one time point at a time and fits the model using all remaining time points. This is useful for very small datasets because it uses as much data as possible for fitting while still checking how well the model predicts data not used for that fit. It is a gentle diagnostic and may be less demanding when adjacent time points are highly related. \\
Leave-contiguous-block-out & leave\_contiguous\_block\_out & Sets aside adjacent time points as a block in x-axis order. The block-size setting controls how many adjacent observations are set aside in each split. This is the primary time-aware sensitivity check when the goal is to test whether the fitted curve can bridge a missing time region. The maximum useful block size depends on the number of observations. Larger blocks may leave too few observations for stable fitting, especially in small datasets. \\
None & none & Disables cross-validation. This may be appropriate when the dataset is too small for meaningful CV, when repeated fitting is unstable, or when only summaries of the fitted model are needed. \\
\bottomrule
\end{longtable}
